% Transcription — fonds Alexandre Grothendieck, shelfmark 12, batch 1 (pages 1-20).
% Established from the facsimile archives/batches/12/batch-01.pdf (a mirror of
% https://grothendieck.umontpellier.fr/12.pdf), per the skill
% transcribe-grothendieck. The batch file carries no cover sheet: PDF page N is
% archivists' page N; the pencilled numbers bottom-left ("2", "3" ... "17",
% "19", "20") read at the PDF page of the same number.
%
% Pass: Opus 5.5 (claude-opus-5-5), 2026-09-23 — first pass, unchecked against the pages by a human.
%
% What the batch is. One continuous draft in his hand on IHES letterhead
% paper (Bures-sur-Yvette), headed « Esquisse d'une théorie des motifs » and
% paginated by him 1 to 17 top right, plus one list on a separate leaf.
% Archivists' pages 1-9 are his 1-9; page 10 is a blank letterhead sheet
% (absent); pages 11-18 are his 10-17 (offset of one). Page 19 is blank
% (absent). Page 20 is a separate leaf (typescript show-through on the other
% face), a boxed « Motifs » heading and a table of 19 numbered topics.
% Nothing is dated on the leaves.
%
% The argument, by his numbered sections:
%   1-5    (p. 1) programme: review of cohomology theories, idea of motives;
%          M^+(k) effective motives, a semi-simple abelian ⊗-category with
%          realisation functors T_DR, T_Hdg, T_B, T^ℓ; grading, h^i(X);
%          M^{+1}(k)° ≃ isogeny category of abelian varieties; Q(-1) and
%          the category M(k) with internal Hom and its universal property.
%   6      (pp. 2-3) level (niveau) of a motive; decomposition by pure level;
%          link with Lefschetz and with the ℓ-adic « filtration géométrique »
%          Filt^j H^i by kernels to open complements of codimension >= j.
%   7      (pp. 3-5) ⊗-categories: rigid, K-linear, connected; Modf(G) as
%          universal example; (C,F) ≃ (Modf(G), oubli), full faithfulness;
%          type fini ⟺ G of finite type, semi-simple ⟺ pro-reductive;
%          twisting fibre functors by torsors, F' = P ×^G F, G' = ad(P);
%          existence of fibre functors over finite extensions, gerbes.
%   8      (pp. 5-6) motives over a regular connected base S; the Betti fibre
%          functor at a complex point gives the motivic Galois group G;
%          Serre's counter-example in char p (supersingular elliptic curves),
%          ℓ-adic fibre functors, G_K over a finite extension K of Q.
%   9      (pp. 6-9) relation with π_1: F associated to a geometric point ξ;
%          examples; lim G_α/G_α^0 ≃ π_1(S, ξ) (boxed (*)), via the
%          π_1-modules of finite rank as motives of level and weight 0 under
%          Tate/Hodge; étale coverings S'/S; (*) bijective iff G connected
%          over algebraically closed fields. Page 7's first version is
%          cancelled by pencil cross-hatching and restarted on page 8.
%   10     (pp. 9, 11-13; his 9-12) relation with profinite π_1: compatibility
%          with Tate's functor, ρ_ℓ : π_1 → G(Q_ℓ), η_ℓ ρ_ℓ = id; a
%          triangle with dashed arrows through G_α/G_α^0; Tate conjecture as
%          G_α = algebraic hull of the ℓ-adic image; finite fields:
%          Frob_{k,F} in the centre Z(Q), G commutative, extension of Ẑ by a
%          (pro)torus; the family of Frobenius f_s in I_G(Q) = T/W.
%   12     (pp. 14-16; his 13-15) supplementary structures: gradings as
%          T = D_K(M) → G, weight y : G_m → G, Hodge i_1, i_2 with
%          i_1 i_2 = y and complex conjugation, ε : G → G_m from Q(-1) with
%          ε(y(λ)) = λ² (boxed), filtrations via Gr(F), torsors under the
%          unipotent H(α). There is no section 11 on the leaves: 10 is
%          followed by 12.
%   13     (pp. 16-18; his 15-17) relation with the Mumford-Tate group; the
%          Hodge conjecture read as G = G_0.
%   14     (p. 18) title « Structure de positivité » only; the draft stops.
%   20     « Motifs »: 19-point plan (M^+(X), variances, T_ℓ, Q_ℓ(n), Hom and
%          RHom, filtrations, rings M^+(X), M(X), Tate, Galois invariants,
%          absolute cohomology, positive forms, L-functions, Hodge).
%
% Readings to note. The centre of G is written with a figure-3 letter, set
% as \mathfrak{Z} (hand.md §4). "½" is his semi- (½-simple, ½-invariants).
% The fast connecting prose often does not carry; the formulas do.
\documentclass[11pt,a4paper]{article}
% Le préambule commun à toutes les transcriptions du fonds Grothendieck.
%
% Il tient en une idée : **l'incertitude se compose, elle ne se résout pas.**
% Une transcription qui lisse un mot douteux a détruit la seule chose pour
% laquelle on la faisait — un lecteur doit pouvoir savoir, à chaque endroit,
% ce qui était sur le feuillet et ce qui a été ajouté. Les six macros qui
% suivent sont donc l'appareil critique, et non de la décoration.
%
% Les mêmes commandes sont reconnues par `scripts/render.mjs`, qui produit la
% vue de lecture du site. Une macro ajoutée ici doit l'être aussi là-bas,
% faute de quoi le rendu échouera bruyamment — ce qui est voulu : un rendu
% silencieusement faux est pire qu'un rendu absent.

\ProvidesPackage{grothendieck}[2026/08/08 Transcriptions du fonds Grothendieck]

\usepackage{amsmath,amssymb,amsthm}
\usepackage{mathrsfs}
\usepackage{stmaryrd}
% Les diagrammes commutatifs sont partout dans ce fonds : c'est la langue
% même de la géométrie algébrique de Grothendieck, et une transcription qui
% les rendrait en prose ne transcrirait rien.
\usepackage{tikz-cd}
% Français par défaut — c'est la langue des feuillets — mais l'anglais est
% chargé aussi : la traduction et le résumé sont des documents anglais, et la
% typographie française (espaces avant les deux-points) y serait une faute.
% Ces documents basculent par \selectlanguage{english} en tête de corps.
\usepackage[english,french]{babel}
\usepackage{fontspec}
\usepackage{geometry}
\geometry{a4paper,margin=25mm,marginparwidth=28mm}
\usepackage{marginnote}
\usepackage{xcolor}
% ulem, not soul. `soul`'s \st and \ul are built on a character-by-character
% scan that XeLaTeX's font machinery breaks: the document fails with an
% undefined control sequence rather than degrading. ulem does the same two jobs
% and is Unicode-engine safe. `normalem` stops it hijacking \emph, which the
% transcriptions use for Grothendieck's own emphasis.
\usepackage[normalem]{ulem}
\usepackage{hyperref}
\hypersetup{colorlinks=true,linkcolor=black,urlcolor=black,pdfborder={0 0 0}}

% --- Métadonnées du lot -----------------------------------------------------
% Reprises telles quelles par le script de rendu, qui les affiche en tête de
% la vue de lecture. Elles fixent ce que le fichier prétend couvrir : sans
% elles, une transcription flotte sans point d'attache dans un fonds de seize
% mille feuillets.

\newcommand{\folder}[1]{\gdef\grothFolder{#1}}
\newcommand{\batch}[1]{\gdef\grothBatch{#1}}
\newcommand{\leaves}[2]{\gdef\grothFirst{#1}\gdef\grothLast{#2}}
\newcommand{\foldertitle}[1]{\gdef\grothFoldertitle{#1}}
\gdef\grothFolder{?}\gdef\grothBatch{?}\gdef\grothFirst{?}\gdef\grothLast{?}\gdef\grothFoldertitle{}

% --- Le résumé --------------------------------------------------------------
%
% Il ouvre la lecture modernisée, sous le titre « Résumé », et il est composé
% en petit. Ce n'est pas un ornement : c'est le seul passage du document qui
% ne soit pas des mathématiques, et un lecteur doit pouvoir le reconnaître
% avant de l'avoir lu — puis le sauter s'il connaît déjà le sujet, ou s'y
% arrêter s'il ne le connaît pas. Le filet qui le suit marque la reprise du
% corps.

\newenvironment{resume}
  {\small\setlength{\parskip}{0.45em}}
  {\par\medskip\hrule\medskip}

% --- Mots-clés ---------------------------------------------------------------
%
% En anglais, en clôture du résumé de la lecture modernisée : le vocabulaire
% moderne sous lequel chercher ce que les feuillets construisent. Le manifeste
% du site (`npm run manifest`) les extrait pour étiqueter la cote entière —
% cette ligne est la source unique des tags, il n'y a pas de fichier à côté.
\newcommand{\keywords}[1]{\par\smallskip\noindent
  {\footnotesize\textsc{Keywords} --- \textit{#1}}\par}

% --- L'appareil critique ----------------------------------------------------

% Le numéro de feuillet, en marge. C'est l'ancre qui permet de régler un
% désaccord sur une formule en regardant *un* feuillet plutôt qu'une chemise
% de six cents. Le site s'en sert aussi pour tourner les pages du fac-similé
% quand on fait défiler la transcription.
\newcommand{\leaf}[1]{%
  \marginnote{\footnotesize\color{gray!70}\textsf{#1}}%
  \label{leaf:#1}\ignorespaces}

% Illisible : on ne devine pas. Un mot inventé qui se lit comme les autres est
% le pire résultat possible.
\newcommand{\ill}{\textcolor{red!60!black}{[\,\ldots\,]}}

% Lecture douteuse : le mot est donné, et signalé comme incertain.
\newcommand{\uncertain}[1]{\uline{#1}}

% Ajout de l'éditeur — une lettre manquante, un mot sous-entendu.
\newcommand{\add}[1]{\textcolor{blue!50!black}{[#1]}}

% Biffé par Grothendieck lui-même. Ce qu'il a rayé fait partie du document :
% une rature dit souvent où la pensée a changé de direction.
\newcommand{\struck}[1]{\sout{#1}}

% Note du transcripteur, sur l'état matériel du feuillet.
\newcommand{\note}[1]{\par\smallskip{\small\color{gray!60!black}\textit{[#1]}}\par\smallskip}

% Note marginale de Grothendieck, distincte du corps du texte.
\newcommand{\marginal}[1]{\par\smallskip\noindent\hspace{1em}%
  {\small\color{gray!60!black}#1}\par\smallskip}

% --- Titre ------------------------------------------------------------------

\AtBeginDocument{%
  \begin{center}
    {\large\bfseries Fonds Alexandre Grothendieck --- cote n\textsuperscript{o}~\grothFolder}\\[2pt]
    {\itshape\grothFoldertitle}\\[4pt]
    {\small feuillets \grothFirst--\grothLast\ (lot \grothBatch)}\\[2pt]
    {\footnotesize\color{gray!60!black}Transcription établie d'après le fac-similé numérisé
     par l'Université de Montpellier.\\
     Les lectures incertaines sont soulignées, les illisibles notées [\,\ldots\,],
     les ajouts entre crochets.}
  \end{center}
  \vspace{1em}}

\endinput


\folder{12}
\batch{1}
\pages{1}{20}
\dating{1964-[vers 1971]}
\watermark{Édition de démonstration}
\foldertitle{Généralités (sous-groupes de Galois motiviques) (1964 ?) : notes manuscrites (s.d.), tapuscrit (s.d.)}

\begin{document}

\section{Esquisse d'une théorie des motifs}
\note{le titre est de sa main, souligné, en tête du feuillet 1 ; la suite est
paginée par lui de 1 à 17, en haut à droite}

\page{1}
1. Revue des théories cohomologiques habituelles pour les schémas projectifs
lisses. L'idée de motifs, \marginal{qui doit \uncertain{remplacer} la
cohomologie entière (la théorie des modules (de t.f.\ sur $\mathbb{Z}$ !))}
\uncertain{Résumé} : a) Isomotifs, b) Motifs ½ simples c) Base ponctuelle.
\note{les parenthèses de cette ligne sont un ajout interlinéaire, d'une encre
plus pâle ; trois astérisques en marge gauche, en face des items 1, du
paragraphe « Les motifs de poids $\geq 1$ » et de « Connexion continue »}

2. Définition de la catégorie $\mathcal{M}^{+}(k)$ des motifs effectifs sur
$k$. C'est une catégorie abélienne « ½ simple » avec $\otimes$. Foncteurs
$T^{*}_{\mathrm{DR}}$, $T_{\mathrm{Hdg}}$ (si $k$ de car.\ nulle)
$T_{B}$ (si $k = \mathbb{C}$) $T^{\ell}$ (si $\ell \neq \operatorname{car} k$)
\struck{\ill{}}. Les foncteurs sont multiplicatifs. Tous ces foncteurs sont
fidèles.
\note{« ½ simple » est interlinéaire, au-dessus de « abélienne » ; après
« car $k$ », un mot biffé puis une parenthèse entièrement noircie}

3. Graduation dans la catégorie $\mathcal{M}^{+}(k)$. (Le produit $\otimes$
compatible avec graduation. Les foncteurs $h^{i}(X)$.
\uncertain{$\mathcal{M}^{+}(X)$}).

4. L'équivalence $\mathcal{M}^{+1}(k)^{\circ} \xrightarrow{\ \sim\ }
\mathrm{Isab}(k)$ ($X \mapsto \mathrm{Pic}^{00}_{X}$,
$A \mapsto h^{1}(\hat{A})$).
\struck{Les motifs $\mathbb{Q}(-1)$, cette catégorie $\mathcal{M}(k)$ des
motifs}
\note{l'exposant de $\mathrm{Pic}$ est lu $00$, sous réserve}

Les motifs de poids $\geq 1$ comme généralisation des VA. La connexion
« discrète » est conservée (s'exprimant ici par le fait que si
$k \to k'$ est une extension de corps algébriquement clos, alors
$\mathcal{M}^{e}(k) \to \mathcal{M}^{e}(k')$ est pleinement fidèle).
\note{l'exposant de $\mathcal{M}$ est douteux : un petit « e », peut-être
autre chose}

Connexion continue… Exemple d'invariant suggéré par la théorie des
\uncertain{schémas} abéliens : caractérisation des motifs \emph{lisses} sur
$S$ \struck{réguliers (\ill{})} \uncertain{normal}.
\note{« normal » est écrit au-dessus du passage biffé}

5. Le motif $\mathbb{Q}(-1)$ (anti-Tate). La catégorie $\mathcal{M}(k)$ des
motifs (isomotifs semi-simples sur $k$ !). Maintenant il y a un
\emph{Hom} interne, en plus du $\otimes$ et de la graduation. Si
$\mathcal{C}$ est une catégorie pseudo-abélienne avec $\otimes$, alors les
foncteurs compatibles avec $\otimes$ de $\mathcal{M}(k)$ dans $\mathcal{C}$
correspondent aux foncteurs $T : \mathcal{V}(k)^{\circ} \to \mathcal{C}$,
satisfaisant Künneth, et tels que \struck{\ill{}}
$T(\mathbb{Q}(-1)) = \operatorname{Im} T(\pi)$ ($\pi$ le projecteur dans
$h(\mathbb{P}^{1}_{k})$ \struck{\ill{}}, d'image
$h^{2}(\mathbb{P}^{1}_{k}) = \mathbb{Q}(-1)$) \emph{soit inversible}. On
montre que les $T$ commutent automatiquement à la formation des \emph{Hom}.
La catégorie des motifs $\mathcal{M}(k)$ est donc caractérisée, avec sa
structure $\otimes$, \struck{\ill{}} à équivalence près par propriété
universelle.

\page{2}
\marginal{\uncertain{Résulte} de la \uncertain{présence} de $\mathcal{M}^{+}$
dans $\mathcal{M}$.}
\note{en tête du feuillet, au-dessus du paragraphe 6, et relié par une
accolade à la définition qui suit}

6. \emph{Niveau d'un motif}. Un motif $M$ est dit \emph{de niveau}
$\leq i$ s'il est de la forme $M'(n)$, où $M' \in \operatorname{Ob}
\mathcal{M}^{+}(k)$ est à composantes homogènes de degrés $\leq i$. On dit
que $M$ est \uncertain{purement} de niveau $i$ s'il est de niveau $\leq i$, et
si aucun sous-motif $\neq 0$ n'est de niveau $\leq i-1$. \marginal{Cet $i$ est
unique, et $\geq 0$.} Tout motif $M$ se décompose de façon unique en somme de
motifs purs de niveau $i$
\[
  M = \coprod_{i \in \mathbb{Z}} M_{i}
\]
donc s'il est de \struck{degré} poids $n$, \marginal{(les $i$ tels que
$M_{i} \neq 0$ ont la même parité que $n$, et \ill{})} \struck{il s'écrit}
\[
  M = \coprod_{i} M'_{i}\Big(\frac{n-i}{2}\Big) \qquad
  M'_{i} \in \mathcal{M}^{+i}_{\text{pur}}(k)
\]
où $M'_{i}$ est \struck{purement de poids} \struck{effectif} pur de poids et
niveau $i$.
\note{« $\mathcal{M}^{+i}_{\text{pur}}(k)$ » est entouré ; « effectif »,
interlinéaire, paraît biffé avec le reste}

\struck{\ill{}} Les motifs de \struck{poids} niveau $\leq i$ sont ceux qui,
modulo twist de Tate, « proviennent » de variétés $X \in \mathcal{V}(k)$ de
dimension $\leq i$. Ceci résulte des faits que $h(X)$ est de niveau
$\leq \dim X$, et que \struck{\uncertain{Addition}} si $i \leq \dim X$,
$h^{i}(X) \hookrightarrow h^{i}(Y)$, où $Y$ est une section linéaire de
dim.\ $i$ (\uncertain{sens} du « th.\ de Lefschetz »). Autre interprétation en
termes de cohomologie $\ell$-adique \struck{(en supposant le corps de base
alg.\ clos, ce qui est licite, car la notion de niveau, \ill{} la \ill{} du
niveau, est invariant par extension du corps de base).} On regarde
$H^{i}(\overline{X}, \mathbb{Q}_{\ell})$ et la « filtration géométrique »
définie par
\[
  \operatorname{Filt}^{j} H^{i}(\overline{X}, \mathbb{Q}_{\ell}) =
  \varinjlim_{U} \operatorname{Ker}\big(H^{i}(\overline{X}, \mathbb{Q}_{\ell})
  \to H^{i}(\overline{U}, \mathbb{Q}_{\ell})\big)
\]
$U$, $U$ ouvert tel que $\operatorname{codim}(X-U) \geq j$.

\marginal{Exemple : si $i \geq n = \dim X$, alors $h^{i}(X)$ est de
\emph{niveau} $\leq 2n-i < i$}
\marginal{Ex.\ Si $i \geq n = \dim X$, $h^{i}(X)$ est de niveau
$\leq 2n-i$}
\note{deux notes marginales à gauche, presque identiques, l'une en face de la
définition, l'autre en face de l'argument de Lefschetz}

\page{3}
Alors on a, \struck{si $X$ provenant de dim.\ $n$ et $i \leq n$}
\uncertain{(supp.\ $i \geq n$)}
\[
  \operatorname{Filt}^{j} H^{i}(\overline{X}, \mathbb{Q}_{\ell}) =
  \text{partie de } T^{\ell}(h^{i}(X))
\]
\struck{provenant} image du plus grand sous-motif \struck{(\ill{}
$2n-i \geq j$)} de niveau $\leq i-2j$ de $h^{i}(X)$.

Cette caractérisation \marginal{du niveau des motifs} est \ill{} \ill{}
\uncertain{courante} si on admet la résolution des singularités
\marginal{des variétés de dim.\ $\leq \dim X = n$ seulement},
\struck{\ill{}} (conséquence de la résolution et des conjectures standard),
on trouve donc que les dimensions des
$\operatorname{Filt}^{j} H^{i}(\overline{X}, \mathbb{Q}_{\ell})$ sont
indépendantes de $\ell$. (Bien mieux, ils correspondent à des motifs, et les
quotients gradués successifs sont \struck{\ill{}} purs de poids $i-2j$
resp.\ $2n-i-j$). Parler des inv.\ birationnels…

\marginal{Conjectures de Tate et de Hodge pourraient venir ici.}
\note{en marge gauche, en oblique, entre le paragraphe précédent et le
paragraphe 7, avec un trait qui les sépare}

7. \emph{Structure des $\otimes$-catégories}. \struck{standard}

\struck{Catégorie} Catégorie abélienne avec opération $\otimes$ « ass.,
comm. » et « unitaire ». On suppose de plus l'existence de \emph{Hom}
internes, au sens du $\otimes$. On dit qu'elle est « rigide » si
\[
  \check{M} \otimes N \xrightarrow{\ \sim\ } \underline{\mathrm{Hom}}(M, N)
\]
Supposons de plus que $\mathcal{C}$ soit $K$-linéaire ($K$ un corps) (i.e.\ on se donne $K \to \operatorname{End}(\mathrm{id}_{\mathcal{C}})$), et que
les \emph{Hom} soient de dim.\ finie sur $k$ : \struck{\ill{}} On dira qu'on
a une $\otimes$-catégorie sur $k$ (qu'on suppose rigide). Une telle catégorie
est dite « connexe » si
$\operatorname{End}(\mathbb{1}_{\mathcal{C}}) \simeq k$ (\uncertain{bien}
l'identité).
\note{il écrit tantôt $K$, tantôt $k$, pour le corps de base de la
$\otimes$-catégorie ; la transcription suit la page}

\page{4}
Exemple : Si $G$ est un groupe algébrique sur $K$, la catégorie
$\mathrm{Modf}(G)$ des \uncertain{modules} sur lesquels $G$ opère.
\uncertain{Plus} \uncertain{avec} groupes pro-algébriques
(\uncertain{p.\ ex.\ profinis}).

Cet exemple \struck{est} essentiellement universel, comme on va voir.

$\otimes$-foncteurs. Foncteurs fibres \marginal{sur $K$, sur $K'$, ….}
Considérons un couple $(\mathcal{C}, F)$ formé d'une $\otimes$-catégorie sur
$K$, avec foncteur fibre rat.\ $/K$. À cette situation est associé
canoniquement un groupe \struck{alg.} affine $G$ sur $K$ (ou pro-groupe
\uncertain{profini} sur $K$), et une équivalence
\[
  (\mathcal{C}, F) \simeq (\mathrm{Modf}(G), \text{oubli}).
\]
De façon précise, associant à chaque $G$ le couple
$((\mathrm{Modf}(G), G), \text{oubli})$, on trouve un foncteur
\[
  \text{groupes affines sur } K \longrightarrow
  \begin{array}{l}
    \otimes\text{-catégories avec foncteur fibre donné, avec} \\ \otimes\text{-foncteurs modulo isom., compatibles avec } F
  \end{array}
\]
et ce foncteur est pleinement fidèle.
$\mathcal{C}$ « de type fini » $\Longleftrightarrow$ $G$ de type fini

$\mathcal{C}$ ½-simple $\Longleftrightarrow$ $G$ pro-réductif.
\struck{\ill{}}
\note{sous la seconde équivalence, « si $K$ de car.\ nulle » ; la fin de la
ligne a été réécrite deux fois, et une ligne entière biffée en dessous}

\marginal{$G(A) = \operatorname{Aut} F_{A}$, ($F_{A}(M) = F(M) \otimes A$)
\struck{\ill{}} comme foncteur fibre à valeurs dans $\mathrm{Mod}(A)$}
\note{note entourée, en marge gauche, en face de l'équivalence
$(\mathcal{C}, F) \simeq (\mathrm{Modf}(G), \text{oubli})$}

\emph{Cor.} Possibilité d'une définition des autres opér.\ tensorielles telles
$\Lambda$, $\mathrm{Symm}$ (en fait, on voit que c'est indépendant de
l'existence et du choix de $F$). Si on a un foncteur fibre, tout autre
foncteur fibre $F'$ s'obtient en « tordant » avec un pro-torseur $P$ sous $G$,
$P = (P_{i})$ :
\[
  F' = P \times^{G} F
\]

\marginal{Remarque : affirmer que si $\mathcal{C}$ de type fini, deux
foncteurs fibres deviennent isomorphes sur ext.\ finie}

\page{5}
Alors $G' = P \times^{G} G = \operatorname{ad}(P)$.

En fait, il y a des cas intéressants où on ne peut trouver aucun foncteur
fibre. Mais si $\mathcal{C}$ est de type fini, alors on
\marginal{doit pouvoir} \uncertain{montrer} qu'on peut trouver un foncteur
fibre sur une extension finie de $K$.
\marginal{(Quand il en est ainsi, \struck{on démontre})}
\struck{On construit ainsi \ill{} ; \ill{}} Admettant le résultat
« d'existence » de foncteurs fibres, on \struck{doit} arrive à décrire
$(\mathcal{C}, F)$ intrinsèquement, en termes d'une gerbe sur un site
convenable…

8. \emph{Cas des catégories de motifs.} \struck{Lien avec le groupe
fondamental profini}

$S$ schéma régulier \marginal{connexe}. On définit $\mathcal{M}^{+}(S)$,
$\mathcal{M}(S)$ en remplaçant ce qui est fait sur un corps, en travaillant
avec $X$ projectifs lisses sur $S$, avec
$C(X) \overset{\mathrm{df}}{=} C(X_{\eta})$ ($\eta$ pt générique de $S$).
Ainsi, $\mathcal{M}(S) \to \mathcal{M}(\eta)$ est pleinement fidèle,
permettant d'interpréter $\mathcal{M}(S)$ comme sous-catégorie pleine de
$\mathcal{M}(\eta)$. De plus, pour tout morphisme $f : S' \to S$ de schémas
réguliers, on définit $f^{*} : \mathcal{M}(S) \to \mathcal{M}(S')$ un
$\otimes$-foncteur correspondant (c'est ici que sert l'hypothèse de
\emph{régularité} sur les schémas….).
\marginal{à expliciter …..}

Tout point \emph{complexe} \marginal{géom.} $\xi$ de $S$ définit un foncteur
fibre « de Betti »
\struck{$F_{\xi}^{\mathrm{Betti}} = F_{\xi}$}
$F_{\xi}^{\mathrm{Betti}} = F_{\xi}^{B} : \mathcal{M}(S) \to
\mathrm{Modf}(\mathbb{Q})$, d'où un groupe pro-algébrique $G$ sur
$\mathbb{Q}$, pro-réductif

\page{6}
en fait, \struck{\ill{}} décrivant $\mathcal{M}(S)$. Pratiquement, on n'a guère
qu'à considérer dans \struck{une} chaque question une sous-catégorie de type
fini de $\mathcal{M}(S)$, soit $\mathcal{C}$. Elle est donc décrite par un
\struck{\ill{}} groupe algébrique \marginal{réductif} $G$ sur $\mathbb{Q}$.
Question intéressante : comment sont reliés ces groupes pour des $\xi$
différents (ils se déduisent les uns des autres par twist…). Déjà intéressant
si $S$ est le spectre d'une ext.\ finie $k$ de $\mathbb{Q}$….!

Si $S$ est de car.\ $p > 0$, on ne peut en général trouver de foncteur fibre
sur $\mathbb{Q}$ (contre-exemple de Serre, lié aux courbes elliptiques de
Hasse nul). Mais on a un foncteur fibre sur $\mathbb{Q}_{\ell}$ pour tout pt
géom.\ de car.\ $\neq \ell$, d'où un foncteur fibre sur une extension finie
$K/\mathbb{Q}$ (qu'on se \uncertain{contentera} : une
$\mathcal{C} \subset \mathcal{M}(S)$ stable par \struck{\ill{}}
$\otimes$, $\vee$, $\oplus$ : engendrée finie…). Donc on trouve un groupe
algébrique réductif $G_{K}$ sur $K$, qui permet de représenter la catégorie
des $K'$-motifs pour tout corps $K'$ \uncertain{qui est une} extension de
$K$.

\marginal{Structure $\varepsilon$, $y$ ($\varepsilon y(\lambda) = \lambda^{2}$),
structure $y$, $i_{1}$, $i_{2}$ \struck{pour corps de base $\mathbb{C}$,
$T_{\mathrm{Hdg}}$}. (A) \ill{}$/G^{0} \to \pi_{1}^{1}$.}
\note{ajout interlinéaire serré, en travers de la ligne du titre 9), qui
annonce les paragraphes 12) et suivants ; le premier symbole de « (A) » et la
parenthèse entre $y$ et « structure » sont lus sous réserve}

9) \emph{Relation avec le $\pi_{1}$ \struck{profini}} : Soit $F$ un foncteur
fibre à coefficients dans $K$ sur $\mathcal{M}(S)$, $\xi$ un pt
géométrique. On dit que $F$ et $\xi$ sont associés si on a un isom.
\note{le numéro 9) paraît récrit sur un 8) ; en marge, quatre petits traits.
La lettre dans « $\mathcal{M}(S)$ » ressemble à un $C$}

\page{7}
\[
  F \mid \mathcal{M}^{+0}(S) \simeq F_{\xi}
\]
où $F_{\xi}$ est le foncteur induit sur $\mathcal{M}^{+0}(S)$ : fibre en
$\xi$. (NB. $S$ étant connexe, ces foncteurs sont isomorphes entre eux, donc
la \emph{condition} d'association ne dépend pas de $\xi$. Mais il faut la
considérer plutôt comme une \emph{donnée} supplémentaire, comme donnée de
l'isom.\ en question). \emph{Ex 1} Si $S = \operatorname{Spec}(\mathbb{C})$,
\struck{alors} $F = F^{\mathrm{Betti}}_{\xi}$, \marginal{\emph{Ex.\ 2}
$T^{\ell}_{\xi}$} alors on a une telle donnée. \emph{Ex 3} Si
$S = \operatorname{Spec} k$, $k$ corps de car.\ nulle (p.\ ex.\ $k = \mathbb{Q}$) et si $F$ est le foncteur fibre \struck{\ill{}}
\marginal{à coefficients dans $k$} $T_{\mathrm{DR}}$ \struck{\ill{}} ou
$T_{\mathrm{Hdg}}$, alors il n'est pas en général associé aux pts
géométriques. \emph{Ex 4} Si \struck{$k$ est} $S$ est le spectre d'un corps
alg.\ clos, alors tout $F$ est associé à $\xi$, du moins si $k$ est
dénombrable \marginal{i.e.\ de degré de transc.\ dénombrable (dénombrable)},
dans $\mathcal{M}(k)$, à engendrement dénombrable….
\note{les numéros des exemples 3 et 4 sont récrits sur d'autres chiffres}

\struck{Alors on trouve un isom.\ canonique
$\pi_{1}(S, \xi) \simeq \varprojlim G_{\alpha}/G_{\alpha}^{0}$ $\{$ les
$G_{\alpha}/G_{\alpha}^{0}$ sont des groupes constants sur $K$ ……. Cela
résulte du \ill{} $\otimes$-foncteur canonique
$\mathrm{Mod}\text{-}\mathrm{cont}(\pi_{1}) \simeq
\mathcal{M}^{+0}(S)$ $\to \mathcal{M}(S)$ qui définit $G \to \pi_{1}$}
\note{tout ce passage est barré au crayon de hachures croisées ; il est repris
en tête du feuillet suivant ; l'indice de « $\mathrm{Mod}$ », lu « cont », est douteux}

\page{8}
Alors on trouve un $\otimes$-foncteur \marginal{pleinement} fidèle
\[
  \pi_{1}\text{-}\mathbb{Q}\text{-modules continus de rang fini sur }
  \mathbb{Q} \longrightarrow \text{motifs sur } S
\]
qui se décrit par un \marginal{épi}morphisme
\[
  G = \varprojlim G_{\alpha} \longrightarrow \pi_{1}(S, \xi)
\]
et se factorise donc en un \emph{isomorphisme} canonique
\[
  (*) \qquad \boxed{\varprojlim_{\alpha} G_{\alpha}/G_{\alpha}^{0}
  \xrightarrow{\ \sim\ } \pi_{1}(S, \xi)}
\]
Soit $\pi_{1}^{!}$ le \ill{} \ill{}. Alors \marginal{supposant $K = \mathbb{Q}$}
\struck{les} (\struck{$\pi_{1}^{!}$} pour simplifier) les
$\pi_{1}^{!}$-$\mathbb{Q}$-modules correspondent à la sous-catégorie de
$\mathcal{M}(S)$ formée des motifs $M$ tels que la sous-$\otimes$-catégorie
engendrée par $M$ soit de « longueur finie ». Utilisant conjecture de Tate
(ou de Hodge si $S$ de car.\ générique nulle), on \struck{voit} prouve que ce
sont les motifs de niveau zéro \marginal{et poids zéro} (i.e.\ les motifs
effectifs de poids zéro), i.e.\ que $(*)$ est un \emph{isomorphisme}. Les
\struck{résultats} réflexions précédentes peuvent se développer aussi
(utilisant la notion de motif sur $S$ à coeff.\ dans $K$) si on ne suppose
pas $K = \mathbb{Q}$.

\struck{\ill{}}

\emph{NB} Soit $S' \to S$ un revêtement étale \marginal{connexe} de $S$,
\marginal{\uncertain{s'envoyant}} et supposons que $F$ provienne d'

\page{9}
un \struck{\ill{}} $\otimes$-foncteur \struck{fibre} sur $S'$. Alors, comme en
théorie des groupes fondamentaux ordinaires, $G' \to G$ est un
monomorphisme, dont l'image est le \struck{noyau} l'image inverse
\marginal{dans $G$} du sous-groupe ouvert de $\pi_{1}(S, \xi)$ qui correspond
à $S'/S$. On en conclut que la bijectivité de $(*)$
\struck{équivalente pour \ill{} $S = \ill{}$} (pour tout $S$, disons) équivaut
au fait que pour $S = \operatorname{Spec} k$, $k$ alg.\ clos, les groupes
$G$ sont \emph{connexes}.

10) \emph{Relation avec le $\pi_{1}$ profini} : \emph{les homs
$\pi_{1} \to G(\mathbb{Q}_{\ell})$}.

Supposons \struck{\ill{}} que $K \hookrightarrow \mathbb{Q}_{\ell}$ ;
\marginal{on dit que} \struck{et que} $F$ \struck{est associé au pt géom.\ $\xi$} \emph{compatible avec le foncteur de Tate} $T^{\ell}_{\xi}$
(\uncertain{i.e.\ on se donne le} plongement donné de $K$ dans
$\mathbb{Q}_{\ell}$) si
\[
  F \otimes_{K} \mathbb{Q}_{\ell} \xrightarrow{\ \sim\ } T^{\ell}_{\xi}
\]
(donnée de compatibilité avec foncteur de Tate). \emph{Exemple 1}
$K = \mathbb{Q}_{\ell}$, $F = T^{\ell}_{\xi}$, \emph{Exemple 2}
$K = \mathbb{Q}$, $F = T^{\xi}_{\mathrm{Betti}}$, où $\xi$ un pt complexe.

Alors \struck{\ill{}} l'homomorphisme canonique
\[
  \pi_{1}(S, \xi) \to \operatorname{Aut} T^{\ell}_{\xi}
\]
\struck{définit} puis :
$\operatorname{Aut} T^{\ell}_{\xi} \simeq
\operatorname{Aut} F \otimes_{K} \mathbb{Q}_{\ell} \simeq G(\mathbb{Q}_{\ell})$
définit un homomorphisme
\[
  \boxed{\rho_{\ell} : \pi_{1}(S, \xi) \longrightarrow G(\mathbb{Q}_{\ell})}
\]

\page{11}
\note{p.\ 10 de sa pagination ; à partir d'ici, son numéro est inférieur d'une
unité à celui des archivistes, jusqu'au feuillet 18 (son 17)}
Noter que si $F$ est compatible avec $T^{\ell}_{\xi}$, il est a fortiori
\emph{associé} à $\xi$, et on trouve
\[
  G \longrightarrow \pi_{1}(S, \xi)
\]
d'où
\[
  \eta_{\ell} : G(\mathbb{Q}_{\ell}) \longrightarrow \pi_{1}(S, \xi) .
\]
On voit tout de suite que cet homomorphisme \struck{est} est inverse à gauche
de $\rho_{\ell}$
\[
  \boxed{\eta_{\ell} \rho_{\ell} = \mathrm{id}_{\pi_{1}(S, \xi)}}
\]
Si on s'intéresse à une sous-catégorie $\mathcal{C}_{\alpha}$ de
$\mathcal{C}$, correspondant à un $G_{\alpha}$, alors
$G_{\alpha}/G_{\alpha}^{0}$ \marginal{définit} \struck{correspond à} un
quotient $\pi_{\alpha}$ de $\pi_{1}(S, \xi)$, et on a :

\begin{tikzcd}
  & G_{\alpha}/G_{\alpha}^{0} \arrow[dr, dashed] & \\
  \pi_{1}(S, \xi) \arrow[r, "\rho_{\ell}"] \arrow[rr, bend right=40, "\text{hom can.}"'] & G_{\alpha}(\mathbb{Q}_{\ell}) \arrow[r] \arrow[u, dashed] & \pi_{\alpha}
\end{tikzcd}

(commutatif).

La conjecture de Tate \marginal{(avec action semi-simple de Galois !)}
équivaut à dire que, \ill{} localisé d'un schéma de type fini sur
$\mathbb{Z}$, $G_{\alpha}$ est l'enveloppe algébrique de
$\operatorname{Im}(\pi_{1}(S, \xi) \to G_{\alpha}(\mathbb{Q}_{\ell}))$.
\struck{NB $G_{\alpha}(\mathbb{Q}_{\ell})$ \ill{}}
\note{à droite, un encadré de quatre ou cinq lignes commençant par « NB »,
barré au crayon de hachures serrées et illisible dessous}

\marginal{\emph{NB} La conjecture de Tate sous la forme : les ½-invariants
dans $H^{2i}(\overline{X}, \mathbb{Q}_{\ell}(i))$ sont les comb.\ à coeff.\ dans $\mathbb{Q}_{\ell}$ des classes algébriques — implique (avec conjectures
standard) la ½ simplicité dans Tate…}

\emph{Cas $S = \operatorname{Spec} k$, $k$ corps fini}.
$\mathcal{C} \subset \mathcal{M}(k)$, et $F$ un « foncteur fibre » sur
$\mathcal{C}$ à coeff.\ dans $K$. Comme on a un élément canonique
\[
  \mathrm{Frob}_{k} \in \operatorname{Aut} \underline{\mathrm{id}}_{\mathcal{M}(k)}
  \longrightarrow \operatorname{Aut} \mathrm{id}_{\mathcal{C}}
  \quad [\simeq \mathfrak{Z}(\mathbb{Q}) \text{, où } \mathfrak{Z}
  \text{ est le centre de } G]
\]
on en déduit un élément canonique de $G(K) = \operatorname{Aut} F$, et même
de $\mathfrak{Z}(\mathbb{Q})$ (où $\mathfrak{Z}$ est le centre de $G$)
\[
  \boxed{\mathrm{Frob}_{k,F} \in \mathfrak{Z}(\mathbb{Q}) \subset G(K)} .
\]
\marginal{N.B. $\mathfrak{Z}$ est \struck{\ill{}} défini sur $\mathbb{Q}$ !}
\note{il note le centre par une lettre en forme de 3, transcrite
$\mathfrak{Z}$ ; dans le crochet, « où $\mathfrak{Z}$ est le centre de $G$ »
est récrit sur un mot biffé}

\page{12}
Si $K \subset \mathbb{Q}_{\ell}$ et si $F$ est compatible avec
$T^{\ell}_{\xi}$, alors l'image de
$\mathrm{Frob}_{k,F} \in \mathfrak{Z}(K) \to \mathfrak{Z}(\mathbb{Q}_{\ell})
\subset G(\mathbb{Q}_{\ell})$ \struck{n'est autre \ill{} par} n'est autre que
l'image du générateur
\[
  \mathrm{Frob}_{k} \in \pi_{1}(k, \xi) \simeq \hat{\mathbb{Z}}
\]
par l'hom.
\[
  \rho_{\ell} : \hat{\mathbb{Z}} \longrightarrow G(\mathbb{Q}_{\ell})
\]
La \emph{conjecture de Tate} pour le corps fini $k$ équivaut à dire que $G$
est l'enveloppe algébrique du sous-groupe engendré par
$\mathrm{Frob}_{k,F} \in G(K)$. Elle implique donc que $G$ est commutatif.
[Comme il est défini \uncertain{pour toute} \ill{}, on voit qu'il est en fait
défini sur $\mathbb{Q}$, canonique, et évidemment on aura en fait, donc
\[
  \mathrm{Frob}_{k,F} \in G(\mathbb{Q})
\]
].
\note{devant $G(\mathbb{Q})$, un $\mathfrak{Z}(\mathbb{Q})$ biffé}

\emph{Conséquence de conjecture de Tate pour $k$}.
\struck{(canoniquement \ill{})} $G$ est extension
\marginal{de $\hat{\mathbb{Z}}$ par un tore resp.\ protore.}

\emph{Conséquences} : Soit maintenant $S$ quelconque, mais muni d'un pt
fermé $s$ tel que $k(s)$ soit fini. Soit $\mathcal{C} \subset \mathcal{M}(S)$,
\struck{\ill{}} \struck{$\mathcal{C}'$ une sous-catégorie de} muni d'un
foncteur fibre $F$, à coeff.\ dans un corps $K \supset \mathbb{Q}$,

\page{13}
d'où un $G$ \uncertain{schéma} en groupes sur $K$. Soit $I_{G}$ le « schéma
des invariants », isomorphe si $G$ connexe à $T/W$ ($T$ tore maximal
\struck{\ill{}}, $W$ le groupe de Weyl). Alors on trouve un élément $f_{s}$
de $I_{G}(\mathbb{Q})$. (NB. $I_{G}$ est en fait défini sur $\mathbb{Q}$).
On trouve ainsi une famille d'éléments $(f_{s})_{s \in S_{0}}$, où $S_{0}$
désigne l'ens.\ des pts $s$ de $S$ tels que $k(s)$ soit fini : famille des
frobenius. Compatibilité avec les homs
\[
  G \longrightarrow \pi_{1}(S, \xi)
\]
et les homs
\[
  \pi_{1}(S, \xi) \longrightarrow G(\mathbb{Q}_{\ell}) .
\]

\page{14}
12) \emph{Structures supplémentaires sur le groupe de Galois motivique} :
\marginal{structures associées aux} graduations, filtrations, au motif de
Tate.

On choisit une sous-\marginal{$\otimes$-}catégorie grad.\ \struck{La
catégorie} $\mathcal{C}$ de $\mathcal{M}(S)$, avec un foncteur fibre $F$,
\marginal{à coeff.\ dans $K$} d'où groupe de Galois motivique $G$.

a) \struck{Une} graduation de type $M$ sur le foncteur fibre, compatible avec
$\otimes$, est équivalente à la donnée d'un hom.
\[
  T = D_{K}(M) \longrightarrow G
\]
La graduation provient d'une graduation fonctorielle de $\mathcal{M}$
\marginal{$\mathrm{df}$ / \ill{}} ssi l'hom $T \to G$ est central. Par suite :

1°) La graduation \marginal{canonique} de type $\mathbb{Z}$ des objets de
$\mathcal{C}$ définit un hom canonique
\[
  y : \mathbb{G}_{m,K} \longrightarrow G
\]

2°) Si $F = T_{\mathrm{Hdg}}$, la bigraduation de ce foncteur correspond : un
hom \struck{$D_{K}(\mathbb{Z}^{2}) = \mathbb{G}_{m}^{2} \to G$ donc :
donnant \ill{}}
\[
  i : \mathbb{G}_{m,K}^{2} \longrightarrow G
\]
donc : deux homs commutant l'un à l'autre
\[
  i_{1}, i_{2} : \mathbb{G}_{m} \longrightarrow G
\]
Comme la graduation totale associée est celle de 1°), on trouve
\[
  i_{1}(\lambda) i_{2}(\lambda) = y(\lambda)
\]
(ce qui contient déjà que $i_{1}$, $i_{2}$ commutent…).
\note{dans $D_{K}(M)$, la lettre sous la parenthèse est surchargée ; lue $M$}

\page{15}
3°) Supposons $S = \operatorname{Spec}(\mathbb{C})$, prenons
$F = T_{\mathrm{Betti}}$ \struck{ou $T_{\mathrm{Hdg}}$}. Alors on
\struck{peut prendre} a $K = \mathbb{Q}$. D'autre part
$F \otimes_{\mathbb{Q}} \mathbb{C} \simeq T_{\mathrm{Hdg}}$
\struck{\ill{}}, donc on a
$G^{\mathrm{Betti}} \otimes_{\mathbb{Q}} \mathbb{C} \simeq G^{\mathrm{Hdg}}$
\struck{\ill{}}. Donc dans $G^{\mathrm{Betti}} \otimes_{\mathbb{Q}} \mathbb{C}$
on trouve \struck{\ill{}} $i_{1}$, $i_{2}$
\[
  i_{1}, i_{2} : \mathbb{G}_{m} \to G^{\mathrm{Betti}} \otimes_{\mathbb{Q}}
  \mathbb{C}
\]
tels que
\[
  i_{1}(\lambda) i_{2}(\lambda) = y(\lambda)
\]
De plus on a
\[
  \overline{i_{1}(\lambda)} = i_{2}(\lambda)
\]
ce qui exprime que les deux graduations \uncertain{relatives} sont conjuguées
complexes l'une de l'autre.
\note{en tête, au-dessus du numéro de page, « $\mathrm{Bett}$ / $\mathrm{Hdg}$ »
biffé}

b) Il y a lieu de supposer que $\mathcal{C} \ni$
\uncertain{$\mathbb{Q}(-1)$}. Alors $\mathbb{Q}(-1)$ \struck{est} définit un
G-module $F(\mathbb{Q}(-1))$ de rang 1, d'où un hom
\[
  \varepsilon : G \longrightarrow \mathbb{G}_{m,K}
\]
Comme $\mathbb{Q}(-1)$ est de \struck{degré} poids 2, on aura
\[
  \boxed{\varepsilon(y(\lambda)) = \lambda^{2}}
\]

c) Supposons que sur le foncteur $F$ il y ait une filtration fonctorielle,
compatible avec gr.\ Comment l'expliciter ? On considère
\[
  F' = \mathrm{Gr}(F)
\]

\page{16}
c'est un autre foncteur fibre, mais gradué, donc défini par un \emph{torseur}
$P$ \struck{\ill{}} sous $G$, et $F$ en termes de $F'$ \ill{} par torseur
$P'$ sous $G'$. Supposons que sur \struck{\ill{}} une extension $K'$ de $K$,
on puisse trouver un isom
\[
  F \simeq F'
\]
compatible avec les filtrations, et induisant l'identité sur les gradués
associés. [\emph{Ex} \marginal{$S = \operatorname{Spec} k$,} $k$ de car.\ nulle, $F = T_{\mathrm{DR}}$ avec filtration canonique ; ça marche aussi si
$S$ est un préschéma \marginal{régulier} de car.\ nulle…]. Alors il y a un
sous-torseur \marginal{canonique} $Q$ de $P$ \struck{\ill{}}, de groupe $H$ le
sous-groupe de $G$ des automorphismes compatibles avec filtration et
induisant l'identité sur le gradué, et de \struck{\ill{}} $Q'$ de $P'$,
\marginal{groupe $H'$ !} On remarque que $H$, $H'$ sont unipotents. Ainsi,
$F$ se déduit canoniquement du foncteur gradué $F'$ [graduation définie par
$\mathbb{G}_{m,K} \xrightarrow{\ \alpha\ } G'$] à l'aide d'un torseur $Q'$
sous le groupe \marginal{unipotent} $H(\alpha)$ de $G'$ défini par $\alpha$
comme on pense.

\marginal{expliciter la construction de $H(\alpha)$}
\note{note entourée en marge gauche}

13) \emph{Relation avec le groupe de \struck{Galois} Mumford-Tate}.

Les structures de Hodge polarisables forment une $\otimes$-catégorie
semi-simple. Toute sous-$\otimes$-catégorie de type fini est « définie » par
un groupe \uncertain{réd.} \ill{} sur $\mathbb{Q}$. Soit $(V, V_{0}, (V^{pq}))$
une structure de Hodge. \ill{}
\note{la dernière ligne, sous un trait horizontal au bas du feuillet, en
partie hors du cadre de la numérisation : « \ill{} avec
$i_{1}, i_{2} : \mathbb{G}_{m,\mathbb{C}} \to G_{\mathbb{C}}$ satisfaisant
$i_{1}(\lambda) i_{2}(\lambda) = i(\lambda)$ central,
$i_{2}(\lambda) = \overline{i_{1}(\lambda)}$ », suivie d'une ligne coupée}

\page{17}
Considérons $\Gamma = \underline{\mathrm{Aut}}_{\mathbb{Q}}(V_{0})$, groupe
algébrique sur $\mathbb{Q}$, et les deux homs
\[
  i_{1}, i_{2} : \mathbb{G}_{m,\mathbb{C}} \longrightarrow \Gamma_{\mathbb{C}}
\]
définissant la bigraduation de $V$
\[
  i(\lambda, \mu) \text{ est } \lambda^{p}\mu^{q} \text{ sur } V^{pq}
\]
\note{au-dessus de la formule, un début biffé, lu $i_{1}(\lambda)$}
Soit $G$ le plus petit sous-groupe algébrique \marginal{de $\Gamma$}
\struck{tel} que $G_{\mathbb{C}}$ contienne $\operatorname{Im} i_{1}$,
$\operatorname{Im} i_{2}$ ; \struck{Alors} c'est aussi le sous-groupe
algébrique de $\Gamma_{\mathbb{C}}$ engendré par les $i_{\sigma}$,
\uncertain{$\sigma$} parcourant les \struck{\ill{}} automorphismes de
$\mathbb{C}$. C'est un groupe algébrique réductif, appelé \emph{groupe de
Mumford-Tate} de la structure de Hodge.
\struck{\ill{} représentation de $G$ \ill{} $M$ \ill{} sur $\mathbb{Q}$,}
\note{un passage de deux lignes et demie, encadré et barré de boucles}
Si $\Sigma$ est une opération tensorielle sur un module, de sorte que $G$
opère sur $\Sigma(V_{0})$, on voit que les éléments de $\Sigma(V_{0})$
invariants sous $G$ \struck{sont \ill{}} sont ceux qui sont invariants par
\marginal{les homs} $i_{1}$, $i_{2}$, i.e.\ qui sont de bidegré $0,0$. Cela
permet de voir que

\page{18}
le « groupe de Galois », au sens structures de Hodge, de
$(V, V_{0}, (V^{pq}))$, s'identifie à $G$.

Supposons maintenant que $V$ provienne d'un motif $M$, engendrant une
\marginal{$\otimes$-}catégorie de motifs $\mathcal{C}_{0}$. On a alors
$\mathcal{C}_{0} \to \mathcal{C}$, donc $G \to G_{0}$, où $G_{0}$ est le
groupe de Galois motivique ($G_{0} \subset \Gamma$). La conjecture de Hodge
s'interprète en disant que \emph{$G = G_{0}$} !

14) \emph{Structure de \uncertain{positivité}}
\note{le paragraphe 14) s'arrête sur son titre ; le reste du feuillet est
vierge. C'est la dernière page de sa pagination 1 à 17}

\section{Motifs}
\note{le titre est de sa main, encadré, en haut à droite du feuillet 20, qui
porte une table en dix-neuf points, sur une feuille dont l'autre face est un
texte dactylographié. Le feuillet 19 est vierge}

\page{20}
\begin{enumerate}
  \item La catégorie $\mathcal{M}^{+}(X)$
  \item Variances pour $X$
  \item Cas $X = \varprojlim X_{i}$
  \item Foncteurs $T_{\ell}$
  \item Les $\mathbb{Q}_{\ell}(n)$
  \item La catégorie $\mathcal{M}(X)$
  \item Les foncteurs $\underline{\mathrm{Hom}}$ et $R\underline{\mathrm{Hom}}$
  \item Motifs constants, tordus, et \uncertain{pt.\ con.}
  \item Filtrations de $\mathcal{M}(X)$ et $\mathcal{M}^{+}(X)$
  \item Anneaux $M^{+}(X)$ et $M(X)$.
  \item Interprétation topologique des \struck{\ill{}} niveaux par la
    filtration canonique \struck{\ill{}} \uncertain{variété proj.\ non
    singulière}.
  \item L'hom $L(K) \to M^{+}(K)$ et invariants birationnels fondamentaux.
    [NB. À généraliser à des coefficients quelconques]
  \item Caractérisation géométrique des filtrations (Conj.\ de Tate)
  \item Invariants de Galois et thms de commutation.
  \item Cohomologie absolue des motifs. Cas des schémas de type fini.
  \item Relation avec les $\mathrm{gr}$ \uncertain{relatifs} et le \ill{} des
    VA sur un schéma de type fini.
  \item Formes positives.
  \item Motifs et fonctions $L$
  \item Relation avec la cohomologie de Hodge….
\end{enumerate}
\note{les numéros 1) à 19) sont les siens ; la liste les porte dans l'ordre}

\end{document}
